\documentclass{article}
\usepackage{amsmath, amssymb, amsthm}
\usepackage{hyperref}
\usepackage{geometry}
\geometry{margin=1in}

\title{Erd\H{o}s Problem \#1068}
\date{Last updated: 2025-10-01}
\author{Source: erdosproblems.com}

\begin{document}
\maketitle

\noindent\textbf{Status:} OPEN \quad \textbf{No prize}

\noindent\textbf{Tags:} graph theory, set theory, chromatic number

\medskip

\noindent\textbf{Problem Statement:}

Does every graph with chromatic number $\aleph_1$ contain a countable subgraph which is infinitely vertex-connected?




This is described in \cite{BoPi24} as a 'version of the Erd\H{o}s-Hajnal problem' (which is [1067] ), but it does not seem to appear in \cite{ErHa66}. We say a graph is infinitely (vertex) connected if any two vertices are connected by infinitely many pairwise vertex-disjoint paths. Soukup \cite{So15} constructed a graph with uncountable chromatic number in which every uncountable set is finitely vertex-connected. A simpler construction was given by Bowler and Pitz \cite{BoPi24}. See also [1067] .




References

[BoPi24] N. Bowler and M. Pitz, A note on uncountably chromatic graphs . arXiv:2402.05984 (2024).

[ErHa66] Erd\H{o}s, P. and Hajnal, A., On chromatic number of graphs and set-systems . Acta Math. Acad. Sci. Hungar. (1966), 61-99.

[So15] Soukup, D\'{a}niel T., Trees, ladders and graphs . J. Combin. Theory Ser. B (2015), 96--116.


\bigskip
\noindent\small{Source: \url{https://www.erdosproblems.com/1068}}
\end{document}
